% Source: https://artofproblemsolving.com/community/c3h2592145_pekkas_700th_post
% Stable identity: aops-topic-2592145-q03
Measure time $t$ in seconds from 12:00:00 and angles in degrees clockwise from 12. The hands are at
\[H=\dfrac{t}{120},\qquad M=\dfrac{t}{10},\qquad S=6t.\]
Two hands are opposite when their angles differ by an odd multiple of $180^\circ$. If two hands are opposite, the third is at $90^\circ$ to both exactly when its angle differs from one of them by an odd multiple of $90^\circ$. Check the three cases.
\begin{itemize}
\item \emph{Hour and minute opposite:} $M-H=\dfrac{11t}{120}=180(2j+1)$, so $t=\dfrac{120\cdot180(2j+1)}{11}$. Then $S-H=\dfrac{719t}{120}=\dfrac{719\cdot180(2j+1)}{11}$. Setting this equal to $90(2m+1)$ gives $1438(2j+1)=11(2m+1)$, with an even left side and an odd right side. Impossible.
\item \emph{Hour and second opposite:} $S-H=\dfrac{719t}{120}=180(2j+1)$. Then $M-H=\dfrac{11t}{120}=\dfrac{11\cdot180(2j+1)}{719}=90(2m+1)$ gives $22(2j+1)=719(2m+1)$: even $=$ odd. Impossible.
\item \emph{Minute and second opposite:} $S-M=\dfrac{59t}{10}=180(2j+1)$, so $t=\dfrac{1800(2j+1)}{59}$. Then $M-H=\dfrac{11t}{120}=\dfrac{165(2j+1)}{59}=90(2m+1)$ gives $11(2j+1)=354(2m+1)$: odd $=$ even. Impossible.
\end{itemize}
So the answer is $\boxed{\text{no}}$: the three hands never form a $180^\circ$ angle and two $90^\circ$ angles. $\blacksquare$
